\section{CONCLUSION}
\label{sec:discussionandconclusion}

This paper presented MotionDAG, a modular runtime framework that targets vendor-agnostic motion pipelines.
Our results show that event-driven, DAG-based motion processing can effectively adapt across the scenarios, integrating diverse XR hardware and personalized humanoid avatars.

We demonstrated three main contributions.
First, our DAG design enables flexible pipeline composition through interface-driven Motion Nodes, allowing users to scale from hand tracking to full-body embodiment without code changes and to adapt pipelines to specific device configurations.
Second, the consistent, asset-driven authoring workflow improved authoring efficiency: participants completed configuration faster, reported lower NASA--TLX workload, and achieved SUS scores in the ``Good'' range, suggesting reduced day-to-day integration burden across devices.
Third, network offloading extends support for resource-constrained devices without modifying core graph semantics.
Because performance gains depend on network conditions and workload characteristics, deployments should re-profile latency percentiles in situ and retain an on-device fallback path to avoid tail-latency degradation under network impairment.
Overall, MotionDAG is well suited to XR deployments spanning heterogeneous XR devices and avatar rigs, where tracking modalities and retargeting pipelines need to be swapped or extended without code refactoring.

Despite its benefits, MotionDAG has limitations that invite further research.
While the framework is designed to provide a consistent interface across platforms, our current evaluation focuses on in-editor authoring UX under a desktop XR simulator setting (excluding build and deploy time) and isolates interface-authoring effort; outcomes may differ in end-to-end HMD deployment workflows that include device-specific calibration and iteration overhead.
Extending evaluation to more complex XR deployments remains an important direction for future work, particularly to assess the impact of the consistent authoring pattern on practical outcomes.

In addition, the current system could be extended with automatic configuration and motion-pipeline presets conditioned on hardware profiles, tracking modalities, and target avatar representations. We expect these extensions to further improve usability.
Algorithmically, our current implementation mainly integrates and extends existing motion-completion methods rather than proposing a new core estimator. The framework would benefit from systematic integration of state-of-the-art pose-estimation approaches, including recent transformer-based and physics-informed models. 
Coordinating heterogeneous methods within a single pipeline remains an open challenge for future development.